\documentclass[11pt]{article}
\usepackage[a4paper,margin=25mm]{geometry}
\usepackage{amsmath,amssymb,amsthm}
\usepackage{longtable,booktabs,array}
\usepackage[hidelinks]{hyperref}
\hypersetup{pdftitle={E37612510\_3: 2x\textasciicircum 4+xy-x+y\textasciicircum 3+y+1=0 has no integer solutions}}
\newtheorem{theorem}{Theorem}
\newtheorem{lemma}[theorem]{Lemma}
% keep the space around binary operators in inline formulas
\medmuskip=4mu plus 2mu\relax
\emergencystretch=1.5em
\tolerance=400
\allowdisplaybreaks
\newcommand{\Z}{\mathbb{Z}}
\newcommand{\Q}{\mathbb{Q}}
\newcommand{\F}{\mathbb{F}}
% Legendre and Jacobi symbols: one fixed size in running text, one in displays
\newcommand{\leg}[2]{\mathchoice{\biggl(\dfrac{#1}{#2}\biggr)}{(#1/#2)}{(#1/#2)}{(#1/#2)}}
\title{\textbf{E37612510\_3}\\[4pt] {\Large The equation $2x^{4}+xy-x+y^{3}+y+1=0$ has no integer solutions}}
\author{}
\date{}
\begin{document}
\maketitle
\vspace{-8ex}
\begin{center}
Size $h=49$, length $l=12$
\end{center}

\begin{theorem}\label{thm:main}
The equation
\begin{equation}\label{eq:main}
2x^{4}+xy-x+y^{3}+y+1=0
\end{equation}
has no integer solutions.
\end{theorem}

The proof uses the cubic Hilbert reciprocity law over the field $K=\Q(\omega)$, where $\omega^2+\omega+1=0$. To a hypothetical integer solution we attach the values of 28 auxiliary polynomials $H_0,\dots,H_{27}$ with coefficients in $\Z[\omega]$, constants $c_0,\dots,c_{27}\in\Z[\omega]$, and a sum $B$ of cubic Hilbert symbols of these values, recorded by a graph with 95 weighted edges. By the reciprocity law, the local contributions $B_p$ of all primes $p$ add up to $0$ modulo $3$. We show that $B_p=0$ for $p\notin S$, compute $B_p$ for the primes $p\in S$, and obtain a contradiction. Two standard facts from class field theory are quoted without proof: the explicit formula for the cubic Hilbert symbol at the places not above $3$, and the Hilbert reciprocity law. Apart from these, the proof uses polynomial identities, which can be checked by expanding both sides, and finite computations with residues.


\section{\texorpdfstring{Cubic Hilbert symbols over $\Q(\omega)$}{Cubic Hilbert symbols over Q(w)}}\label{sec:symbols}

\paragraph{The field.} Let $K=\Q(\omega)$, where $\omega^2+\omega+1=0$. Its ring of integers is $\Z[\omega]$, and every element of $K$ can be written uniquely as $A+B\omega$ with $A,B\in\Q$. We write
\[
\langle A,B\rangle=A+B\omega,
\]
also when $A$ and $B$ are polynomials in $x$ and $y$. Then $\langle A,B\rangle\langle C,D\rangle=\langle AC-BD,\ AD+BC-BD\rangle$, and the norm is $N(\langle A,B\rangle)=A^2-AB+B^2$. In particular, a value $\langle A,B\rangle$ with $A,B\in\Z$ is zero if and only if $A=B=0$.

\paragraph{Places.} The finite places of $K$ correspond to the non-zero prime ideals of $\Z[\omega]$. A prime $p\equiv2\pmod3$ is inert: there is one place above $p$, with uniformizer $p$ and residue field $\Z[\omega]/p\Z[\omega]$ with $p^2$ elements. If $p\equiv1\pmod3$, the polynomial $t^2+t+1$ has two roots $r$ modulo $p$, and for each of them there is a place $v_r$ above $p$, at which $\omega$ is the root of $t^2+t+1$ in $\Z_p$ congruent to $r$ modulo $p$; the completion is $\Q_p$, the uniformizer is $p$, the residue field is $\F_p$, and $\omega$ reduces to $r$. The prime $3$ is ramified: $\pi=1-\omega$ generates the unique prime ideal above $3$, $\pi^2=-3\omega$, and $\Z[\omega]/\pi\Z[\omega]=\F_3$, where $\langle A,B\rangle$ reduces to $A+B$. The only infinite place of $K$ is complex. For a finite place $v$ we write $K_v$ for the completion and $v(a)$ for the normalized valuation, and we call $a$ a \emph{unit} at $v$ if $v(a)=0$.

\paragraph{Hilbert symbols.} For a place $v$ of $K$ and $a,b\in K_v^\times$, the cubic Hilbert symbol $(a,b)_v$ is a cube root of unity, see \cite[Chapter~V, \S3]{N}. We write $(a,b)_v=\omega^{\beta_v(a,b)}$ with $\beta_v(a,b)\in\Z/3\Z$, and we use the following three facts \cite[Chapter~V, \S3 and Chapter~VI, \S8]{N}.
\begin{itemize}
\item[(H1)] $\beta_v(a,bc)=\beta_v(a,b)+\beta_v(a,c)$, $\beta_v(a,b)=-\beta_v(b,a)$, and $\beta_v(a,c^3)=0$. In particular, $\beta_v(a,b)$ depends only on the classes of $a$ and $b$ in $K_v^\times/K_v^{\times3}$.
\item[(H2)] Let $v$ be a finite place not above $3$, let $q$ be the number of elements of its residue field, and let $\varpi$ be a uniformizer at $v$. Write $a=\varpi^eu$ and $b=\varpi^fw$ with units $u$ and $w$, and define $\chi(u)\in\Z/3\Z$ by $\bar u^{(q-1)/3}=\bar\omega^{\chi(u)}$, where the bar denotes the reduction to the residue field. Then $\beta_v(a,b)=e\,\chi(w)-f\,\chi(u)$. In particular, $\beta_v(a,b)=0$ if $a$ and $b$ are both units.
\item[(H3)] (Hilbert reciprocity.) For $a,b\in K^\times$, $\beta_v(a,b)=0$ for all but finitely many places $v$, and $\sum_v\beta_v(a,b)=0$, where the sum is over all places of $K$. At the complex place, $\beta_v(a,b)=0$.
\end{itemize}
The general formula in (H2) contains an additional term $ef\chi(-1)$, which vanishes because $-1=(-1)^3$ is a cube. Formula (H2) fixes the normalization of the symbols; with the opposite sign convention all values $\beta_v$ change sign, and the proof below is the same.

\paragraph{Classes at the places not above $3$.} Let $p\neq3$ be a prime and $v$ a place above $p$. For $a=p^eu\in K_v^\times$ with a unit $u$, put $c_v(a)=(e\bmod3,\ \chi(u))\in\F_3^2$. By (H2) and (H1), $c_v$ is additive, $c_v(a)$ determines the class of $a$ modulo cubes, and $\beta_v(a,b)=e\chi(w)-f\chi(u)$ for $c_v(a)=(e,\chi(u))$ and $c_v(b)=(f,\chi(w))$. We put $c_p(a)=c_v(a)\in\F_3^2$ if $p\equiv2\pmod3$, and $c_p(a)=(c_{v_r}(a),c_{v_{r'}}(a))\in\F_3^4$ if $p\equiv1\pmod3$, where $r<r'$ are the roots of $t^2+t+1$ in $\{1,\dots,p-1\}$; accordingly, $\beta_p$ is the sum of the forms $\beta_v$ over the places $v$ above $p$. For $p\equiv2\pmod3$ and a rational integer $u$ prime to $p$, we have $\bar u\in\F_p$, so that $\bar u^{(p^2-1)/3}=(\bar u^{p-1})^{(p+1)/3}=1$ and $\chi(u)=0$.

\paragraph{The place above $3$.}
\begin{lemma}\label{lem:M3}
Let $v$ be the place above $3$. Every non-zero element of $K_v$ is equal to $\pi^a\omega^b2^c(1+3\omega)^d$ times a cube, for unique $a,b,c,d\in\Z/3\Z$, and every unit congruent to $1$ modulo $9$ is a cube. Put $c_3(A)=(a,b,c,d)\in\F_3^4$. Then $c_3$ is additive, and $\beta_v(A,B)=c_3(A)^{\mathsf T}M_3\,c_3(B)$, where
\[
M_3=\begin{pmatrix}0&0&2&0\\0&0&1&2\\1&2&0&0\\0&1&0&0\end{pmatrix}\pmod3.
\]
\end{lemma}

\begin{proof}
For $\alpha\in\Z[\omega]$, the polynomial $g(s)=s+3s^2+3s^3-\alpha$ satisfies $g'(s)\equiv1\pmod\pi$, so by Hensel's lemma it has a root $s$ in the completion of $\Z[\omega]$, and then $(1+3s)^3=1+9(s+3s^2+3s^3)=1+9\alpha$. Hence every unit congruent to $1$ modulo $9$ is a cube in $K_v$, and the class of a unit modulo cubes depends only on its residue modulo $9$. There are $54$ units in $\Z[\omega]/(9)$, and a direct enumeration shows that the $27$ products $\omega^b2^c(1+3\omega)^d$ with $b,c,d\in\{0,1,2\}$ are pairwise distinct modulo $9$, also after multiplication by $-1=(-1)^3$. Hence they represent all $27$ classes of units modulo cubes. Since $\pi^3$ is a cube times a unit, the exponent of $\pi$ is the valuation modulo $3$, and the first statement follows.

By (H1), it remains to compute $\beta_v$ on the $16$ ordered pairs of the elements $\pi$, $\omega$, $2$ and $1+3\omega$. Their norms are $3$, $1$, $4$ and $7$, hence, for each such pair, all the symbols at the places not above $2$, $3$ and $7$ vanish by (H2), and, by (H3), $\beta_v$ at the place above $3$ is minus the sum of the values at the place above $2$ and at the two places above $7$. These places correspond to $\bar\omega=2$ and $\bar\omega=4$ in $\F_7$. The classes $(e,\chi)$ of the four elements at these places are
\[
\begin{array}{c|ccc}
 & 2 & \bar\omega=2 & \bar\omega=4\\\hline
\pi & (0,2) & (0,0) & (0,2)\\
\omega & (0,1) & (0,2) & (0,2)\\
2 & (1,0) & (0,2) & (0,1)\\
1+3\omega & (0,2) & (1,0) & (0,0)
\end{array}
\]
where, at the places above $7$, $\chi(u)$ is defined by $\bar u^2=\bar\omega^{\chi(u)}$ in $\F_7$, and $7$ is a uniformizer. Formula (H2) and the sum over these three places give $M_3$.
\end{proof}

We write $\beta_3$ for the form of Lemma~\ref{lem:M3}. Since $3=-\omega^2\pi^2$ and $4=2^2$, we have $c_3(3)=(2,2,0,0)$ and $c_3(4)=(0,0,2,0)$.

\paragraph{Residue classes.} The local computations use the following description of the classes of the values of a polynomial on a residue class.

\begin{lemma}\label{lem:residue}
Let $p$ be a prime, $k\geq1$, and let $z\in\Z[\omega]$ be non-zero with $z\equiv\zeta\pmod{p^k}$, where $\zeta=\langle A,B\rangle$ with $0\leq A,B<p^k$. Then $c_p(z)\in b+D$ for the vector $b$ and the subspace $D$ given as follows.
\begin{itemize}
\item[(i)] $p\equiv2\pmod3$. If $\zeta\neq0$, then $b=c_p(\zeta)$ and $D=0$. If $\zeta=0$, then $b=0$, and $D$ is $\F_3^2$, or the span of $(1,0)$ if $z\in\Z$.
\item[(ii)] $p\equiv1\pmod3$. For each root $r$, let $\rho$ be the root of $t^2+t+1$ modulo $p^k$ congruent to $r$. If $A+B\rho\not\equiv0\pmod{p^k}$, the two coordinates of $c_{v_r}(z)$ are those of $c_{v_r}(\zeta)$; otherwise these two coordinates are arbitrary.
\item[(iii)] $p=3$, $z\notin\Z$. If $\zeta\neq0$, let $a=v(\zeta)<2k$ be its valuation at $\pi$ and $m=2k-a$. Then $b=c_3(\zeta)$, and $D=0$ if $m\geq4$, while $D=D_m$ if $m\leq3$, where $D_1$ is spanned by $(0,1,0,0)$, $(0,0,1,0)$ and $(0,0,0,1)$, $D_2$ by $(0,0,1,0)$ and $(0,0,0,1)$, and $D_3$ by $(0,0,1,1)$. If $\zeta=0$, then $b=0$ and $D=\F_3^4$.
\item[(iv)] $p=3$, $z\in\Z$, so that $B=0$. If $3^k\nmid A$, let $j=v_3(A)$. Then $b=c_3(A)$ and $D=0$ if $k-j\geq2$, while $b=j\,c_3(3)$ and $D$ is spanned by $c_3(4)$ if $k-j=1$. If $3^k\mid A$, then $b=0$ and $D$ is spanned by $c_3(3)$ and $c_3(4)$.
\end{itemize}
\end{lemma}

\begin{proof}
(i) If $\zeta\neq0$, then $e=v(z)=v(\zeta)<k$ and $z/p^e\equiv\zeta/p^e\pmod p$, so that $\chi$ is determined. If $z\in\Z$, its unit part is a rational integer, and $\chi=0$. (ii) At $v_r$, $z$ and $\zeta$ are the $p$-adic numbers $A'+B'\omega$ and $A+B\omega$ with $\omega\equiv\rho\pmod{p^k}$, and the argument of (i) applies. (iii) If $\zeta\neq0$, then $v(z)=a$, and $z/\pi^a\equiv\zeta/\pi^a\pmod{\pi^m}$. By Lemma~\ref{lem:M3}, the class of a unit depends only on its residue modulo $9=\omega\pi^4$; the units congruent to $1$ modulo $\pi^m$ form a subgroup, and a direct enumeration of their residues modulo $9$ shows that their classes form $D_m$. (iv) If $3^k\nmid A$, then $z=3^ju$ with a rational unit $u\equiv A/3^j\pmod{3^{k-j}}$. The class of a rational unit depends only on its residue modulo $9$, the units $\pm1$ are cubes, and the rational units congruent to $1$ modulo $3$ are $4^s$ times a unit congruent to $1$ modulo $9$; hence the classes are as stated, and if $3^k\mid A$ the valuation $j$ is also unknown.
\end{proof}

\begin{lemma}\label{lem:envelope}
Let $p$ be a prime, and consider vertices $i$ with vectors $b_i$, subspaces $D_i$ and vectors $\gamma_i$, and a set of edges $(i,j)$ with $i<j$ and weights $w_{ij}\in\{1,2\}$. Put $\rho_{ij}=w_{ij}$ and $\rho_{ji}=-w_{ij}$ for the edges, and $\sigma_i=\gamma_i+\sum_j\rho_{ij}b_j$. Suppose that $\beta_p(d,\sigma_i)=0$ for every $i$ and every $d\in D_i$, and that $\beta_p(d,d')=0$ for every edge $(i,j)$ and all $d\in D_i$, $d'\in D_j$. Then
\[
\sum_{(i,j)}w_{ij}\beta_p(c_i,c_j)+\sum_i\beta_p(c_i,\gamma_i)=\sum_{(i,j)}w_{ij}\beta_p(b_i,b_j)+\sum_i\beta_p(b_i,\gamma_i)
\]
whenever $c_i\in b_i+D_i$ for every $i$. By bilinearity, it suffices to check the hypotheses for $d$ and $d'$ in spanning sets of $D_i$ and $D_j$.
\end{lemma}

\begin{proof}
Write $c_i=b_i+d_i$ with $d_i\in D_i$. As $\beta_p$ is bilinear and $\beta_p(u,v)=-\beta_p(v,u)$, the left side equals the right side plus $\sum_i\beta_p(d_i,\sigma_i)+\sum_{(i,j)}w_{ij}\beta_p(d_i,d_j)$, and both sums vanish.
\end{proof}
\section{The graph}

Put $F(x,y)=2x^{4}+xy-x+y^{3}+y+1$, so that \eqref{eq:main} is the equation $F(x,y)=0$. The certificate consists of 28 polynomials $H_i$ with coefficients in $\Z[\omega]$ and constants $c_i\in\Z[\omega]$, listed in Table~\ref{tab:vert1}, and of 95 \emph{edges} $(i,j)$ with $i<j$ and weights $w_{ij}\in\{1,2\}$, listed below. For an edge $(i,j)$ we put $\rho_{ij}=w_{ij}$ and $\rho_{ji}=-w_{ij}\bmod 3$, and $\rho_{ij}=0$ if $i$ and $j$ are not joined by an edge.
The edges $(i,j;w_{ij})$ are $(0,8;1)$, $(0,9;2)$, $(0,10;2)$, $(0,13;1)$, $(0,20;1)$, $(0,26;2)$, $(0,27;2)$, $(1,4;2)$, $(1,11;2)$, $(1,13;1)$, $(1,17;2)$, $(1,20;2)$, $(1,21;1)$, $(1,22;2)$, $(2,17;1)$, $(2,21;2)$, $(3,7;1)$, $(3,9;1)$, $(3,10;1)$, $(3,14;2)$, $(3,17;1)$, $(3,18;2)$, $(3,24;1)$, $(4,6;2)$, $(4,11;1)$, $(4,22;1)$, $(4,24;2)$, $(4,26;1)$, $(4,27;2)$, $(5,7;2)$, $(5,9;2)$, $(5,12;2)$, $(5,13;2)$, $(5,14;1)$, $(5,19;1)$, $(6,17;1)$, $(6,19;2)$, $(6,23;1)$, $(6,25;2)$, $(7,11;1)$, $(7,13;2)$, $(7,14;2)$, $(7,17;1)$, $(7,19;2)$, $(7,20;1)$, $(7,22;1)$, $(8,13;2)$, $(8,16;2)$, $(8,17;2)$, $(8,18;2)$, $(8,19;2)$, $(8,20;1)$, $(8,23;2)$, $(8,24;2)$, $(8,27;2)$, $(9,14;1)$, $(9,15;1)$, $(9,18;2)$, $(9,26;2)$, $(9,27;2)$, $(10,18;1)$, $(10,20;2)$, $(10,21;2)$, $(11,13;2)$, $(11,17;2)$, $(12,18;1)$, $(12,19;1)$, $(12,21;2)$, $(12,23;2)$, $(13,25;2)$, $(14,21;2)$, $(14,24;2)$, $(14,25;2)$, $(15,16;2)$, $(15,21;1)$, $(15,22;1)$, $(15,24;2)$, $(15,26;2)$, $(15,27;1)$, $(16,19;2)$, $(16,25;2)$, $(16,27;2)$, $(17,20;1)$, $(17,25;1)$, $(18,20;2)$, $(18,21;1)$, $(19,20;2)$, $(20,22;2)$, $(20,27;1)$, $(21,22;1)$, $(21,23;2)$, $(21,25;1)$, $(23,24;1)$, $(24,25;2)$, $(25,26;1)$.
Let $(\xi,\eta)$ denote $(x,y)$ or $(y,x)$. We call a vertex $i$ \emph{linear} if $H_i=\varepsilon_i(\eta-r_i(\xi))$ for a unit $\varepsilon_i$ of $\Z[\omega]$ and a polynomial $r_i$ with coefficients in $\Z[\omega]$; the last column of Table~\ref{tab:vert1} gives $\eta=r_i(\xi)$ for the linear vertices. For a linear vertex $i$, we put $f_i(t)=F$ and $g_{ij}(t)=H_j$ evaluated at $\xi=t$, $\eta=r_i(t)$; these are polynomials in $t$ with coefficients in $\Z[\omega]$.
For every linear vertex $i$, Table~\ref{tab:star} gives an integer $D_i\geq1$ and a polynomial $w_i(t)$ with coefficients in $\Z[\omega]$ such that
\begin{equation}\label{eq:starlin}
D_i^3\,c_i\prod_jg_{ij}^{\rho_{ij}}-w_i^3=f_iQ_i
\end{equation}
for a polynomial $Q_i(t)$ with coefficients in $\Z[\omega]$, where $\rho_{ij}\in\{0,1,2\}$; $Q_i$ is the quotient of the division of the left side by $f_i$, and the remainder is zero.
For every edge $(i,j)$ with a linear endpoint, Table~\ref{tab:sep} gives the first linear endpoint $k\in\{i,j\}$ and the factorization of the norm of the resultant $R_{ij}=\operatorname{Res}_t(f_k,g_{kl})$, where $\{k,l\}=\{i,j\}$; it is non-zero.
These identities can be checked by expanding both sides and by polynomial division, using $\omega^2=-1-\omega$.

\section{Proof of Theorem~\ref{thm:main}}

Suppose that $(x,y)\in\Z^2$ is a solution of \eqref{eq:main}, so that $F(x,y)=0$; from now on the polynomials denote their values at $(x,y)$, which lie in $\Z[\omega]$.
\paragraph{Step 1: the values are non-zero.} For $i\in\{0,\allowbreak 1,\allowbreak 2,\allowbreak 4,\allowbreak 5,\allowbreak 6,\allowbreak 8,\allowbreak 9,\allowbreak 11,\allowbreak 13,\allowbreak 14,\allowbreak 16,\allowbreak 17,\allowbreak 18,\allowbreak 19,\allowbreak 22,\allowbreak 23,\allowbreak 24,\allowbreak 27\}$, there is no residue pair $(a,b)$ modulo $2$ with $F(a,b)\equiv0$ and $H_i(a,b)\equiv0\pmod{2}$, as one checks for the $2^2$ residue pairs; hence $H_i\neq0$. For $i\in\{3,\allowbreak 7,\allowbreak 12,\allowbreak 21,\allowbreak 25,\allowbreak 26\}$, there is no residue pair $(a,b)$ modulo $3$ with $F(a,b)\equiv0$ and $H_i(a,b)\equiv0\pmod{3}$, as one checks for the $3^2$ residue pairs; hence $H_i\neq0$. For $i\in\{20\}$, there is no residue pair $(a,b)$ modulo $5$ with $F(a,b)\equiv0$ and $H_i(a,b)\equiv0\pmod{5}$, as one checks for the $5^2$ residue pairs; hence $H_i\neq0$. For $i\in\{10\}$, there is no residue pair $(a,b)$ modulo $7$ with $F(a,b)\equiv0$ and $H_i(a,b)\equiv0\pmod{7}$, as one checks for the $7^2$ residue pairs; hence $H_i\neq0$. For $i\in\{15\}$, there is no residue pair $(a,b)$ modulo $11$ with $F(a,b)\equiv0$ and $H_i(a,b)\equiv0\pmod{11}$, as one checks for the $11^2$ residue pairs; hence $H_i\neq0$. Also, all $c_i$ are non-zero.

\paragraph{Step 2: reciprocity.} For every place $v$ of $K$, put
\begin{equation}\label{eq:B}
B_v=\sum_{(i,j)}w_{ij}\,\beta_v(H_i,H_j)+\sum_i\beta_v(H_i,c_i)\in\Z/3\Z,
\end{equation}
the first sum over the edges, and, for a rational prime $p$, put $B_p=\sum_{v\mid p}B_v$. Applying (H3) to every pair $(H_i,H_j)$ with an edge, multiplied by $w_{ij}$, and to every pair $(H_i,c_i)$, and adding, we find that $B_v=0$ for all but finitely many $v$ and $\sum_pB_p=0$, as the complex place contributes $0$. By (H1) and the definition of $c_p$, $B_p=\sum_{(i,j)}w_{ij}\beta_p(c_p(H_i),c_p(H_j))+\sum_i\beta_p(c_p(H_i),c_p(c_i))$, where $\beta_p$ is the form of Lemma~\ref{lem:M3} if $p=3$.

\paragraph{Step 3: the primes outside $S$.} Let $S=\{2,3,5,7,11,13\}$. It contains $3$ and the primes dividing the norms of the constants $c_i$, of the $D_i$, of the $R_{ij}$. Let $p\notin S$, and let $v$ be a place above $p$, with residue field $k_v$; then all these numbers are units at $v$. First, no edge $(i,j)$ has both values divisible by $v$. Indeed, suppose that $v$ divides $H_i$ and $H_j$. If $R_{ij}$ is defined, let $k$ be the linear endpoint used for it and $\{k,l\}=\{i,j\}$, and put $t=\xi$. As $H_k=\varepsilon_k(\eta-r_k(\xi))$, we have $\eta\equiv r_k(t)$ modulo $v$, hence $f_k(t)\equiv F(x,y)=0$ and $g_{kl}(t)\equiv H_l\equiv0$ modulo $v$. The resultant can be written as $R_{ij}=Af_k+Cg_{kl}$ with polynomials $A$ and $C$ with coefficients in $\Z[\omega]$, so $v$ divides $R_{ij}$, a contradiction. By (H2), a term of \eqref{eq:B} in which all arguments are units vanishes. Now let $v$ divide $H_i$. Collecting the terms of \eqref{eq:B} which contain $H_i$, and using (H1) and $\rho_{ji}=-w_{ij}$ for the edges $(j,i)$ with $j<i$, we obtain $\beta_v(H_i,\Pi_i)$ with $\Pi_i=c_i\prod_jH_j^{\rho_{ij}}$, which is a unit. We claim that $D_i^3\Pi_i\equiv W^3$ modulo $v$ for some $W\in\Z[\omega]$. If $i$ is linear, put $t=\xi$; then $g_{ij}(t)\equiv H_j$ and $f_i(t)\equiv0$ modulo $v$ as above, and \eqref{eq:starlin} gives the claim with $W=w_i(t)$. Hence the reduction of $\Pi_i$ is a non-zero cube in $k_v$, $\chi(\Pi_i)=0$, and $\beta_v(H_i,\Pi_i)=v(H_i)\chi(\Pi_i)=0$ by (H2). Hence $B_v=0$, and $B_p=0$.

\paragraph{Step 4: the primes in $S$.} Let $p\in S$. Here $2$ is inert, $5$ is inert, the places above $7$ are $v_{2}$ and $v_{4}$, $11$ is inert, and the places above $13$ are $v_{3}$ and $v_{9}$. Suppose that $x\equiv a$ and $y\equiv b\pmod{p^k}$. Then $H_i\equiv H_i(a,b)\pmod{p^k}$, and Lemma~\ref{lem:residue}, applied with $\zeta$ the residue of $H_i(a,b)$, gives a vector $b_i$ and a subspace $D_i$ with $c_p(H_i)\in b_i+D_i$; when $H_i$ has rational coefficients, $H_i\in\Z$, and we use this in (i) and (iv). With $\gamma_i=c_p(c_i)$, if the hypotheses of Lemma~\ref{lem:envelope} hold, then $B_p$ is determined by $a$, $b$ and $k$. Starting from the solutions of \eqref{eq:main} modulo $p$, we refine a residue class $x\equiv a$, $y\equiv b\pmod{p^k}$ to the classes modulo $p^{k+1}$ that contain solutions of \eqref{eq:main} modulo $p^{k+1}$, until these hypotheses hold, or until no solution remains. Table~\ref{tab:local} lists the resulting terminal classes and the value of $B_p$ in each; every integer solution lies in one of them. The result is $B_{2}=0$, $B_{3}=1$, $B_{5}=0$, $B_{7}=0$, $B_{11}=0$, $B_{13}\in\{0,1\}$.

\paragraph{Step 5: contradiction.} By Steps~2 and~3, $\sum_{p\in S}B_p=0$. But by Step~4, $\sum_{p\in S}B_p\in\{1,2\}$ in $\Z/3\Z$. This contradiction proves the theorem.
\qed


\begin{thebibliography}{9}
\bibitem{N} J.~Neukirch, \emph{Algebraic Number Theory}, Grundlehren der mathematischen Wissenschaften 322, Springer, Berlin, 1999.
\end{thebibliography}

\clearpage
\begingroup\small
\setlength{\LTcapwidth}{\textwidth}
\begin{longtable}{@{}rlll@{}}
\caption{The polynomials $H_i$, the constants $c_i$ and, for the linear vertices, $\eta=r_i(\xi)$. Here $\langle A,B\rangle=A+B\omega$.}\label{tab:vert1}\\
\hline $i$ & $H_i$ & $c_i$ & chart\\\hline\endfirsthead\hline $i$ & $H_i$ & $c_i$ & chart\\\hline\endhead\hline\endfoot
0 & $\langle -y+1,1\rangle$ & $\langle 1,0\rangle$ & $y=\langle 1,1\rangle$\\
1 & $\langle -y-1,0\rangle$ & $\langle -4,-4\rangle$ & $y=\langle -1,0\rangle$\\
2 & $\langle -y-1,-1\rangle$ & $\langle -4,-8\rangle$ & $y=\langle -1,-1\rangle$\\
3 & $\langle 2x-y,x+1\rangle$ & $\langle 0,-3\rangle$ & $y=\langle 2x,x+1\rangle$\\
4 & $\langle x-y-1,x^{2}\rangle$ & $\langle -1,-1\rangle$ & $y=\langle x-1,x^{2}\rangle$\\
5 & $\langle x^{2}+x-y-1,x\rangle$ & $\langle -15,-7\rangle$ & $y=\langle x^{2}+x-1,x\rangle$\\
6 & $\langle -x+1,1\rangle$ & $\langle 1,0\rangle$ & $x=\langle 1,1\rangle$\\
7 & $\langle -x+1,0\rangle$ & $\langle -2,-2\rangle$ & $x=\langle 1,0\rangle$\\
8 & $\langle -x,1\rangle$ & $\langle 1,0\rangle$ & $x=\langle 0,1\rangle$\\
9 & $\langle -x,-1\rangle$ & $\langle 0,-1\rangle$ & $x=\langle 0,-1\rangle$\\
10 & $\langle -x-1,0\rangle$ & $\langle 1,-1\rangle$ & $x=\langle -1,0\rangle$\\
11 & $\langle -x,-y\rangle$ & $\langle -2,-1\rangle$ & $y=\langle x,x\rangle$\\
12 & $\langle -x+y+1,y\rangle$ & $\langle 0,-6\rangle$ & $y=\langle 0,-x+1\rangle$\\
13 & $\langle -x,y+1\rangle$ & $\langle 1,0\rangle$ & $y=\langle -x-1,-x\rangle$\\
14 & $\langle -x+y,y+1\rangle$ & $\langle -3,-3\rangle$ & $y=\langle -1,-x-1\rangle$\\
15 & $\langle -x+y-1,0\rangle$ & $\langle 2,-2\rangle$ & $y=\langle x+1,0\rangle$\\
16 & $\langle -x-1,y-1\rangle$ & $\langle -2,-2\rangle$ & $y=\langle -x,-x-1\rangle$\\
17 & $\langle -y,-x\rangle$ & $\langle -3,-3\rangle$ & $y=\langle 0,-x\rangle$\\
18 & $\langle x-y,0\rangle$ & $\langle 1,0\rangle$ & $y=\langle x,0\rangle$\\
19 & $\langle -x-1,-y-1\rangle$ & $\langle -12,0\rangle$ & $y=\langle x,x+1\rangle$\\
20 & $\langle 2x^{2}-y,0\rangle$ & $\langle 0,-12\rangle$ & $y=\langle 2x^{2},0\rangle$\\
21 & $\langle -x+1,-y\rangle$ & $\langle 1,0\rangle$ & $y=\langle x-1,x-1\rangle$\\
22 & $\langle -x^{2}-y,0\rangle$ & $\langle -2,-4\rangle$ & $y=\langle -x^{2},0\rangle$\\
23 & $\langle x^{2}-y,0\rangle$ & $\langle 4,0\rangle$ & $y=\langle x^{2},0\rangle$\\
24 & $\langle -y,1\rangle$ & $\langle 4,0\rangle$ & $y=\langle 0,1\rangle$\\
25 & $\langle -y,-x^{2}-x\rangle$ & $\langle 1,0\rangle$ & $y=\langle 0,-x^{2}-x\rangle$\\
26 & $\langle -x-y-1,-y\rangle$ & $\langle -12,-12\rangle$ & $y=\langle 0,x+1\rangle$\\
27 & $\langle -x^{2}-y,-x^{2}-x\rangle$ & $\langle -1,-1\rangle$ & $y=\langle -x^{2},-x^{2}-x\rangle$\\
\end{longtable}
\endgroup

\begingroup\scriptsize
\setlength{\LTcapwidth}{\textwidth}
\begin{longtable}{@{}ll@{}}
\caption{The data of the identities \eqref{eq:starlin}; the entry $j{:}e$ means $\rho_{ij}=e\neq0$.}\label{tab:star}\\
\hline\endfirsthead\hline\endhead\hline\endfoot
\multicolumn{2}{@{}p{0.97\textwidth}@{}}{$i=0$, $D_i=2$, $\rho_{ij}$: $8{:}1, \allowbreak 9{:}2, \allowbreak 10{:}2, \allowbreak 13{:}1, \allowbreak 20{:}1, \allowbreak 26{:}2, \allowbreak 27{:}2$}\\
$w_i$ & $\langle 2t^{3}+2t^{2}+t+1,$\\
 & $\quad 6t^{3}+6t^{2}+t-2\rangle$\\
\addlinespace[2pt]
\multicolumn{2}{@{}p{0.97\textwidth}@{}}{$i=1$, $D_i=2$, $\rho_{ij}$: $4{:}2, \allowbreak 11{:}2, \allowbreak 13{:}1, \allowbreak 17{:}2, \allowbreak 20{:}2, \allowbreak 21{:}1, \allowbreak 22{:}2$}\\
$w_i$ & $\langle -4t^{3}+2t^{2}+2t+2,$\\
 & $\quad -4t^{3}-2t\rangle$\\
\addlinespace[2pt]
\multicolumn{2}{@{}p{0.97\textwidth}@{}}{$i=2$, $D_i=1$, $\rho_{ij}$: $17{:}1, \allowbreak 21{:}2$}\\
$w_i$ & $\langle 0,$\\
 & $\quad -2t^{2}\rangle$\\
\addlinespace[2pt]
\multicolumn{2}{@{}p{0.97\textwidth}@{}}{$i=3$, $D_i=2$, $\rho_{ij}$: $7{:}1, \allowbreak 9{:}1, \allowbreak 10{:}1, \allowbreak 14{:}2, \allowbreak 17{:}1, \allowbreak 18{:}2, \allowbreak 24{:}1$}\\
$w_i$ & $\langle 4t^{3}-8t^{2}-2t+6,$\\
 & $\quad 8t^{3}+2t^{2}-10t\rangle$\\
\addlinespace[2pt]
\multicolumn{2}{@{}p{0.97\textwidth}@{}}{$i=4$, $D_i=1$, $\rho_{ij}$: $1{:}1, \allowbreak 6{:}2, \allowbreak 11{:}1, \allowbreak 22{:}1, \allowbreak 24{:}2, \allowbreak 26{:}1, \allowbreak 27{:}2$}\\
$w_i$ & $\langle 4t^{4}+3t^{2}+t,$\\
 & $\quad -t^{5}-2t^{3}+5t^{2}+1\rangle$\\
\addlinespace[2pt]
\multicolumn{2}{@{}p{0.97\textwidth}@{}}{$i=5$, $D_i=2$, $\rho_{ij}$: $7{:}2, \allowbreak 9{:}2, \allowbreak 12{:}2, \allowbreak 13{:}2, \allowbreak 14{:}1, \allowbreak 19{:}1$}\\
$w_i$ & $\langle 8t^{5}+14t^{4}-32t^{3}+10t^{2}+12t-4,$\\
 & $\quad 26t^{4}-14t^{3}-18t^{2}+18t-2\rangle$\\
\addlinespace[2pt]
\multicolumn{2}{@{}p{0.97\textwidth}@{}}{$i=6$, $D_i=2$, $\rho_{ij}$: $4{:}1, \allowbreak 17{:}1, \allowbreak 19{:}2, \allowbreak 23{:}1, \allowbreak 25{:}2$}\\
$w_i$ & $\langle -2t^{2}-2t+8,$\\
 & $\quad -2t^{2}+4t+4\rangle$\\
\addlinespace[2pt]
\multicolumn{2}{@{}p{0.97\textwidth}@{}}{$i=7$, $D_i=2$, $\rho_{ij}$: $3{:}2, \allowbreak 5{:}1, \allowbreak 11{:}1, \allowbreak 13{:}2, \allowbreak 14{:}2, \allowbreak 17{:}1, \allowbreak 19{:}2, \allowbreak 20{:}1, \allowbreak 22{:}1$}\\
$w_i$ & $\langle 2t^{2}+4t+12,$\\
 & $\quad 4t^{2}+8t-4\rangle$\\
\addlinespace[2pt]
\multicolumn{2}{@{}p{0.97\textwidth}@{}}{$i=8$, $D_i=4$, $\rho_{ij}$: $0{:}2, \allowbreak 13{:}2, \allowbreak 16{:}2, \allowbreak 17{:}2, \allowbreak 18{:}2, \allowbreak 19{:}2, \allowbreak 20{:}1, \allowbreak 23{:}2, \allowbreak 24{:}2, \allowbreak 27{:}2$}\\
$w_i$ & $\langle 44t+40,$\\
 & $\quad 32t^{2}+60t+20\rangle$\\
\addlinespace[2pt]
\multicolumn{2}{@{}p{0.97\textwidth}@{}}{$i=9$, $D_i=4$, $\rho_{ij}$: $0{:}1, \allowbreak 3{:}2, \allowbreak 5{:}1, \allowbreak 14{:}1, \allowbreak 15{:}1, \allowbreak 18{:}2, \allowbreak 26{:}2, \allowbreak 27{:}2$}\\
$w_i$ & $\langle 4t^{2}+8t,$\\
 & $\quad 4t^{2}-4t+8\rangle$\\
\addlinespace[2pt]
\multicolumn{2}{@{}p{0.97\textwidth}@{}}{$i=10$, $D_i=2$, $\rho_{ij}$: $0{:}1, \allowbreak 3{:}2, \allowbreak 18{:}1, \allowbreak 20{:}2, \allowbreak 21{:}2$}\\
$w_i$ & $\langle -4t^{2}+16,$\\
 & $\quad -2t^{2}+8\rangle$\\
\addlinespace[2pt]
\multicolumn{2}{@{}p{0.97\textwidth}@{}}{$i=11$, $D_i=2$, $\rho_{ij}$: $1{:}1, \allowbreak 4{:}2, \allowbreak 7{:}2, \allowbreak 13{:}2, \allowbreak 17{:}2$}\\
$w_i$ & $\langle 2t^{3}-6t+4,$\\
 & $\quad 4t^{3}-6t^{2}+2\rangle$\\
\addlinespace[2pt]
\multicolumn{2}{@{}p{0.97\textwidth}@{}}{$i=12$, $D_i=2$, $\rho_{ij}$: $5{:}1, \allowbreak 18{:}1, \allowbreak 19{:}1, \allowbreak 21{:}2, \allowbreak 23{:}2$}\\
$w_i$ & $\langle -6t^{2}+12t-6,$\\
 & $\quad 6t^{3}-12t^{2}+12t-6\rangle$\\
\addlinespace[2pt]
\multicolumn{2}{@{}p{0.97\textwidth}@{}}{$i=13$, $D_i=2$, $\rho_{ij}$: $0{:}2, \allowbreak 1{:}2, \allowbreak 5{:}1, \allowbreak 7{:}1, \allowbreak 8{:}1, \allowbreak 11{:}1, \allowbreak 25{:}2$}\\
$w_i$ & $\langle -8t^{2}-6t,$\\
 & $\quad 3t^{3}-t^{2}+4t+1\rangle$\\
\addlinespace[2pt]
\multicolumn{2}{@{}p{0.97\textwidth}@{}}{$i=14$, $D_i=4$, $\rho_{ij}$: $3{:}1, \allowbreak 5{:}2, \allowbreak 7{:}1, \allowbreak 9{:}2, \allowbreak 21{:}2, \allowbreak 24{:}2, \allowbreak 25{:}2$}\\
$w_i$ & $\langle 27t^{3}+6t^{2}-3t+3,$\\
 & $\quad -6t^{3}+12t^{2}+15t+15\rangle$\\
\addlinespace[2pt]
\multicolumn{2}{@{}p{0.97\textwidth}@{}}{$i=15$, $D_i=2$, $\rho_{ij}$: $9{:}2, \allowbreak 16{:}2, \allowbreak 21{:}1, \allowbreak 22{:}1, \allowbreak 24{:}2, \allowbreak 26{:}2, \allowbreak 27{:}1$}\\
$w_i$ & $\langle 13t^{3}+8t^{2}-2t-9,$\\
 & $\quad 2t^{3}+16t^{2}+20t+6\rangle$\\
\addlinespace[2pt]
\multicolumn{2}{@{}p{0.97\textwidth}@{}}{$i=16$, $D_i=2$, $\rho_{ij}$: $8{:}1, \allowbreak 15{:}1, \allowbreak 19{:}2, \allowbreak 25{:}2, \allowbreak 27{:}2$}\\
$w_i$ & $\langle -4t^{3}+4t^{2}-10t+2,$\\
 & $\quad -10t^{3}+4t^{2}-2t-2\rangle$\\
\addlinespace[2pt]
\multicolumn{2}{@{}p{0.97\textwidth}@{}}{$i=17$, $D_i=4$, $\rho_{ij}$: $1{:}1, \allowbreak 2{:}2, \allowbreak 3{:}2, \allowbreak 6{:}2, \allowbreak 7{:}2, \allowbreak 8{:}1, \allowbreak 11{:}1, \allowbreak 20{:}1, \allowbreak 25{:}1$}\\
$w_i$ & $\langle -6t^{3}-12t^{2}+6t+2,$\\
 & $\quad -6t^{2}-6t+4\rangle$\\
\addlinespace[2pt]
\multicolumn{2}{@{}p{0.97\textwidth}@{}}{$i=18$, $D_i=2$, $\rho_{ij}$: $3{:}1, \allowbreak 8{:}1, \allowbreak 9{:}1, \allowbreak 10{:}2, \allowbreak 12{:}2, \allowbreak 20{:}2, \allowbreak 21{:}1$}\\
$w_i$ & $\langle -4t^{3}-2t,$\\
 & $\quad 0\rangle$\\
\addlinespace[2pt]
\multicolumn{2}{@{}p{0.97\textwidth}@{}}{$i=19$, $D_i=1$, $\rho_{ij}$: $5{:}2, \allowbreak 6{:}1, \allowbreak 7{:}1, \allowbreak 8{:}1, \allowbreak 12{:}2, \allowbreak 16{:}1, \allowbreak 20{:}2$}\\
$w_i$ & $\langle 3t^{3}-6t^{2}+6t,$\\
 & $\quad 6t^{3}-6t^{2}-3t-3\rangle$\\
\addlinespace[2pt]
\multicolumn{2}{@{}p{0.97\textwidth}@{}}{$i=20$, $D_i=8$, $\rho_{ij}$: $0{:}2, \allowbreak 1{:}1, \allowbreak 7{:}2, \allowbreak 8{:}2, \allowbreak 10{:}1, \allowbreak 17{:}2, \allowbreak 18{:}1, \allowbreak 19{:}1, \allowbreak 22{:}2, \allowbreak 27{:}1$}\\
$w_i$ & $\langle 48t^{5}-72t^{3}+24t^{2}+24t-24,$\\
 & $\quad 12t^{5}+30t^{4}-6t^{3}-36t^{2}+69t-27\rangle$\\
\addlinespace[2pt]
\multicolumn{2}{@{}p{0.97\textwidth}@{}}{$i=21$, $D_i=8$, $\rho_{ij}$: $1{:}2, \allowbreak 2{:}1, \allowbreak 10{:}1, \allowbreak 12{:}1, \allowbreak 14{:}1, \allowbreak 15{:}2, \allowbreak 18{:}2, \allowbreak 22{:}1, \allowbreak 23{:}2, \allowbreak 25{:}1$}\\
$w_i$ & $\langle -11t^{3}-7t^{2}+40t-24,$\\
 & $\quad -19t^{3}-38t^{2}+50t-15\rangle$\\
\addlinespace[2pt]
\multicolumn{2}{@{}p{0.97\textwidth}@{}}{$i=22$, $D_i=22$, $\rho_{ij}$: $1{:}1, \allowbreak 4{:}2, \allowbreak 7{:}2, \allowbreak 15{:}2, \allowbreak 20{:}1, \allowbreak 21{:}2$}\\
$w_i$ & $\langle 60t^{5}-40t^{4}-20t^{3}-44t^{2}+16t+20,$\\
 & $\quad 60t^{5}+4t^{4}-64t^{3}-88t^{2}-28t+64\rangle$\\
\addlinespace[2pt]
\multicolumn{2}{@{}p{0.97\textwidth}@{}}{$i=23$, $D_i=1$, $\rho_{ij}$: $6{:}2, \allowbreak 8{:}1, \allowbreak 12{:}1, \allowbreak 21{:}1, \allowbreak 24{:}1$}\\
$w_i$ & $\langle -2t^{2}+2t,$\\
 & $\quad 2t\rangle$\\
\addlinespace[2pt]
\multicolumn{2}{@{}p{0.97\textwidth}@{}}{$i=24$, $D_i=2$, $\rho_{ij}$: $3{:}2, \allowbreak 4{:}1, \allowbreak 8{:}1, \allowbreak 14{:}1, \allowbreak 15{:}1, \allowbreak 23{:}2, \allowbreak 25{:}2$}\\
$w_i$ & $\langle -6t^{3}+2t^{2}+8t+10,$\\
 & $\quad 6t^{3}+10t^{2}+10t+2\rangle$\\
\addlinespace[2pt]
\multicolumn{2}{@{}p{0.97\textwidth}@{}}{$i=25$, $D_i=676$, $\rho_{ij}$: $6{:}1, \allowbreak 13{:}1, \allowbreak 14{:}1, \allowbreak 16{:}1, \allowbreak 17{:}2, \allowbreak 21{:}2, \allowbreak 24{:}1, \allowbreak 26{:}1$}\\
$w_i$ & $\langle -1352t^{5}-676t^{4}-1352t^{3}-676t^{2}-676t+676,$\\
 & $\quad -676t^{5}+676t^{4}-676t^{3}-676t^{2}-1352t\rangle$\\
\addlinespace[2pt]
\multicolumn{2}{@{}p{0.97\textwidth}@{}}{$i=26$, $D_i=1$, $\rho_{ij}$: $0{:}1, \allowbreak 4{:}2, \allowbreak 9{:}1, \allowbreak 15{:}1, \allowbreak 25{:}2$}\\
$w_i$ & $\langle -t^{3}-3t^{2}-2t-4,$\\
 & $\quad 4t^{3}-3t^{2}-4t-5\rangle$\\
\addlinespace[2pt]
\multicolumn{2}{@{}p{0.97\textwidth}@{}}{$i=27$, $D_i=1$, $\rho_{ij}$: $0{:}1, \allowbreak 4{:}1, \allowbreak 8{:}1, \allowbreak 9{:}1, \allowbreak 15{:}2, \allowbreak 16{:}1, \allowbreak 20{:}2$}\\
$w_i$ & $\langle 3t^{5}+3t^{4}+t^{3}+t^{2},$\\
 & $\quad t^{5}+2t^{4}-t^{3}+t^{2}\rangle$\\
\addlinespace[2pt]
\end{longtable}
\endgroup

\begingroup\small
\setlength{\LTcapwidth}{\textwidth}
\begin{longtable}{@{}lrl@{\qquad}lrl@{}}
\caption{The edges $(i,j)$ with a linear endpoint, the linear endpoint $k$ used, and the norm of $R_{ij}=\operatorname{Res}_t(f_k,g_{kl})$.}\label{tab:sep}\\
\hline $(i,j)$ & $k$ & $N(R_{ij})$ & $(i,j)$ & $k$ & $N(R_{ij})$\\\hline\endfirsthead\hline $(i,j)$ & $k$ & $N(R_{ij})$ & $(i,j)$ & $k$ & $N(R_{ij})$\\\hline\endhead\hline\endfoot
$(0,8)$ & 0 & $2^{2}$ & $(8,17)$ & 8 & $2^{2}$\\
$(0,9)$ & 0 & $2^{2}\cdot 3$ & $(8,18)$ & 8 & $1$\\
$(0,10)$ & 0 & $3^{2}$ & $(8,19)$ & 8 & $2^{2}\cdot 7$\\
$(0,13)$ & 0 & $7^{3}$ & $(8,20)$ & 8 & $7\cdot 13$\\
$(0,20)$ & 0 & $2^{4}\cdot 3\cdot 5^{2}$ & $(8,23)$ & 8 & $2^{2}$\\
$(0,26)$ & 0 & $1$ & $(8,24)$ & 8 & $1$\\
$(0,27)$ & 0 & $3\cdot 5^{2}$ & $(8,27)$ & 8 & $13$\\
$(1,4)$ & 1 & $3\cdot 7$ & $(9,14)$ & 9 & $3\cdot 7$\\
$(1,11)$ & 1 & $1$ & $(9,15)$ & 9 & $2^{2}\cdot 7$\\
$(1,13)$ & 1 & $1$ & $(9,18)$ & 9 & $3$\\
$(1,17)$ & 1 & $1$ & $(9,26)$ & 9 & $1$\\
$(1,20)$ & 1 & $2^{4}\cdot 3^{4}$ & $(9,27)$ & 9 & $3\cdot 7$\\
$(1,21)$ & 1 & $3\cdot 7$ & $(10,18)$ & 10 & $3^{2}$\\
$(1,22)$ & 1 & $3^{2}$ & $(10,20)$ & 10 & $2^{4}\cdot 3^{2}$\\
$(2,17)$ & 2 & $2^{2}$ & $(10,21)$ & 10 & $2^{4}\cdot 3^{2}$\\
$(2,21)$ & 2 & $3$ & $(11,13)$ & 11 & $2^{2}\cdot 7^{2}$\\
$(3,7)$ & 3 & $2^{2}\cdot 7$ & $(11,17)$ & 11 & $3^{4}$\\
$(3,9)$ & 3 & $7$ & $(12,18)$ & 12 & $3$\\
$(3,10)$ & 3 & $2^{4}$ & $(12,19)$ & 12 & $3^{5}$\\
$(3,14)$ & 3 & $2^{2}\cdot 3\cdot 7$ & $(12,21)$ & 12 & $2^{2}\cdot 3^{4}$\\
$(3,17)$ & 3 & $2^{2}\cdot 7^{2}$ & $(12,23)$ & 12 & $2^{4}$\\
$(3,18)$ & 3 & $1$ & $(13,25)$ & 13 & $7^{3}$\\
$(3,24)$ & 3 & $3^{5}$ & $(14,21)$ & 14 & $3^{5}$\\
$(4,6)$ & 4 & $3\cdot 7$ & $(14,24)$ & 14 & $2^{4}\cdot 3\cdot 7$\\
$(4,11)$ & 4 & $1$ & $(14,25)$ & 14 & $2^{2}\cdot 3^{2}$\\
$(4,22)$ & 4 & $2^{4}\cdot 3$ & $(15,16)$ & 15 & $2^{2}\cdot 7^{2}$\\
$(4,24)$ & 4 & $7$ & $(15,21)$ & 15 & $2^{2}\cdot 3\cdot 7$\\
$(4,26)$ & 4 & $2^{8}$ & $(15,22)$ & 15 & $2^{4}$\\
$(4,27)$ & 4 & $3\cdot 7^{2}$ & $(15,24)$ & 15 & $2^{4}\cdot 3\cdot 7$\\
$(5,7)$ & 5 & $7$ & $(15,26)$ & 15 & $2^{4}\cdot 3^{4}$\\
$(5,9)$ & 5 & $2^{2}$ & $(15,27)$ & 15 & $7^{3}$\\
$(5,12)$ & 5 & $3^{2}$ & $(16,19)$ & 16 & $2^{8}$\\
$(5,13)$ & 5 & $7^{2}$ & $(16,25)$ & 16 & $7$\\
$(5,14)$ & 5 & $3^{6}$ & $(16,27)$ & 16 & $7$\\
$(5,19)$ & 5 & $2^{2}\cdot 3^{2}$ & $(17,20)$ & 17 & $2^{2}\cdot 7$\\
$(6,17)$ & 6 & $7$ & $(17,25)$ & 17 & $1$\\
$(6,19)$ & 6 & $3\cdot 7$ & $(18,20)$ & 18 & $2^{6}\cdot 3^{2}$\\
$(6,23)$ & 6 & $2^{2}$ & $(18,21)$ & 18 & $3$\\
$(6,25)$ & 6 & $2^{2}\cdot 7$ & $(19,20)$ & 19 & $2^{2}\cdot 3^{2}\cdot 7^{3}$\\
$(7,11)$ & 7 & $7$ & $(20,22)$ & 20 & $3^{12}$\\
$(7,13)$ & 7 & $7^{2}$ & $(20,27)$ & 20 & $3^{3}\cdot 5^{2}\cdot 7\cdot 13$\\
$(7,14)$ & 7 & $7$ & $(21,22)$ & 21 & $2^{4}\cdot 3$\\
$(7,17)$ & 7 & $7$ & $(21,23)$ & 21 & $2^{4}$\\
$(7,19)$ & 7 & $7$ & $(21,25)$ & 21 & $2^{2}\cdot 3^{2}\cdot 7^{2}$\\
$(7,20)$ & 7 & $2^{2}\cdot 7^{2}$ & $(23,24)$ & 23 & $2^{4}$\\
$(7,22)$ & 7 & $1$ & $(24,25)$ & 24 & $2^{2}$\\
$(8,13)$ & 8 & $1$ & $(25,26)$ & 25 & $2^{8}$\\
$(8,16)$ & 8 & $2^{4}$ & & &\\
\end{longtable}
\endgroup

\begingroup\small
\setlength{\LTcapwidth}{\textwidth}
\begin{longtable}{@{}rrl>{\raggedright\arraybackslash}p{0.66\textwidth}@{}}
\caption{The terminal residue classes $x\equiv a$, $y\equiv b\pmod{p^k}$ of Step~4, grouped by $p$, $k$ and the value of $B_p$.}\label{tab:local}\\
\hline $p$ & $k$ & $B_p$ & classes $(a,b)$\\\hline\endfirsthead
\hline $p$ & $k$ & $B_p$ & classes $(a,b)$\\\hline\endhead\hline\endfoot
2 & 3 & $0$ & $(3,0)$,\allowbreak\ $(3,2)$,\allowbreak\ $(3,4)$,\allowbreak\ $(3,6)$\\
\addlinespace[3pt]
3 & 3 & $1$ & $(0,7)$,\allowbreak\ $(3,16)$,\allowbreak\ $(6,25)$,\allowbreak\ $(9,7)$,\allowbreak\ $(11,5)$,\allowbreak\ $(11,14)$,\allowbreak\ $(11,23)$,\allowbreak\ $(12,16)$,\allowbreak\ $(15,25)$,\allowbreak\ $(18,7)$,\allowbreak\ $(20,2)$,\allowbreak\ $(20,11)$,\allowbreak\ $(20,20)$,\allowbreak\ $(21,16)$,\allowbreak\ $(24,25)$\\
3 & 4 & $1$ & $(2,8)$,\allowbreak\ $(2,35)$,\allowbreak\ $(2,62)$,\allowbreak\ $(29,26)$,\allowbreak\ $(29,53)$,\allowbreak\ $(29,80)$,\allowbreak\ $(56,17)$,\allowbreak\ $(56,44)$,\allowbreak\ $(56,71)$\\
\addlinespace[3pt]
5 & 1 & $0$ & $(1,1)$,\allowbreak\ $(1,3)$,\allowbreak\ $(2,1)$,\allowbreak\ $(2,2)$,\allowbreak\ $(3,0)$,\allowbreak\ $(3,1)$,\allowbreak\ $(3,4)$,\allowbreak\ $(4,1)$\\
\addlinespace[3pt]
7 & 1 & $0$ & $(5,0)$\\
7 & 2 & $0$ & $(2,36)$,\allowbreak\ $(5,36)$,\allowbreak\ $(5,41)$,\allowbreak\ $(9,43)$,\allowbreak\ $(12,15)$,\allowbreak\ $(12,27)$,\allowbreak\ $(16,1)$,\allowbreak\ $(19,13)$,\allowbreak\ $(19,43)$,\allowbreak\ $(23,8)$,\allowbreak\ $(26,22)$,\allowbreak\ $(26,48)$,\allowbreak\ $(30,15)$,\allowbreak\ $(33,1)$,\allowbreak\ $(33,34)$,\allowbreak\ $(37,22)$,\allowbreak\ $(40,20)$,\allowbreak\ $(40,29)$,\allowbreak\ $(43,2)$,\allowbreak\ $(43,9)$,\allowbreak\ $(43,16)$,\allowbreak\ $(43,23)$,\allowbreak\ $(43,30)$,\allowbreak\ $(43,37)$,\allowbreak\ $(44,29)$,\allowbreak\ $(47,6)$,\allowbreak\ $(47,8)$\\
7 & 3 & $0$ & $(1,213)$,\allowbreak\ $(15,171)$,\allowbreak\ $(22,297)$,\allowbreak\ $(36,157)$,\allowbreak\ $(43,234)$,\allowbreak\ $(50,66)$,\allowbreak\ $(57,339)$,\allowbreak\ $(64,24)$,\allowbreak\ $(71,150)$,\allowbreak\ $(85,10)$,\allowbreak\ $(92,87)$,\allowbreak\ $(99,262)$,\allowbreak\ $(113,220)$,\allowbreak\ $(120,3)$,\allowbreak\ $(127,227)$,\allowbreak\ $(134,206)$,\allowbreak\ $(141,283)$,\allowbreak\ $(148,115)$,\allowbreak\ $(162,73)$,\allowbreak\ $(169,199)$,\allowbreak\ $(176,80)$,\allowbreak\ $(183,59)$,\allowbreak\ $(190,136)$,\allowbreak\ $(197,311)$,\allowbreak\ $(204,241)$,\allowbreak\ $(211,269)$,\allowbreak\ $(218,52)$,\allowbreak\ $(232,255)$,\allowbreak\ $(239,332)$,\allowbreak\ $(246,164)$,\allowbreak\ $(260,122)$,\allowbreak\ $(267,248)$,\allowbreak\ $(281,108)$,\allowbreak\ $(288,44)$,\allowbreak\ $(288,93)$,\allowbreak\ $(288,142)$,\allowbreak\ $(288,185)$,\allowbreak\ $(288,191)$,\allowbreak\ $(288,240)$,\allowbreak\ $(288,289)$,\allowbreak\ $(288,338)$,\allowbreak\ $(295,17)$,\allowbreak\ $(309,318)$,\allowbreak\ $(316,101)$,\allowbreak\ $(330,304)$,\allowbreak\ $(337,38)$\\
7 & 4 & $0$ & $(8,486)$,\allowbreak\ $(29,1893)$,\allowbreak\ $(78,1403)$,\allowbreak\ $(106,192)$,\allowbreak\ $(155,45)$,\allowbreak\ $(225,2334)$,\allowbreak\ $(253,2152)$,\allowbreak\ $(274,1844)$,\allowbreak\ $(302,2005)$,\allowbreak\ $(323,1354)$,\allowbreak\ $(351,1858)$,\allowbreak\ $(372,864)$,\allowbreak\ $(421,374)$,\allowbreak\ $(449,1564)$,\allowbreak\ $(498,1417)$,\allowbreak\ $(568,1305)$,\allowbreak\ $(596,1123)$,\allowbreak\ $(617,815)$,\allowbreak\ $(645,976)$,\allowbreak\ $(666,325)$,\allowbreak\ $(694,829)$,\allowbreak\ $(715,2236)$,\allowbreak\ $(764,1746)$,\allowbreak\ $(792,535)$,\allowbreak\ $(841,388)$,\allowbreak\ $(911,276)$,\allowbreak\ $(939,94)$,\allowbreak\ $(960,2187)$,\allowbreak\ $(988,2348)$,\allowbreak\ $(1009,1697)$,\allowbreak\ $(1037,2201)$,\allowbreak\ $(1058,1207)$,\allowbreak\ $(1107,717)$,\allowbreak\ $(1135,1907)$,\allowbreak\ $(1184,1760)$,\allowbreak\ $(1254,1648)$,\allowbreak\ $(1282,1466)$,\allowbreak\ $(1303,1158)$,\allowbreak\ $(1331,1319)$,\allowbreak\ $(1352,668)$,\allowbreak\ $(1380,1172)$,\allowbreak\ $(1401,178)$,\allowbreak\ $(1450,2089)$,\allowbreak\ $(1478,878)$,\allowbreak\ $(1527,731)$,\allowbreak\ $(1597,619)$,\allowbreak\ $(1625,437)$,\allowbreak\ $(1646,129)$,\allowbreak\ $(1674,290)$,\allowbreak\ $(1695,2040)$,\allowbreak\ $(1723,143)$,\allowbreak\ $(1744,1550)$,\allowbreak\ $(1793,1060)$,\allowbreak\ $(1821,2250)$,\allowbreak\ $(1870,2103)$,\allowbreak\ $(1940,1991)$,\allowbreak\ $(1968,1809)$,\allowbreak\ $(1989,1501)$,\allowbreak\ $(2017,1662)$,\allowbreak\ $(2038,1011)$,\allowbreak\ $(2066,1515)$,\allowbreak\ $(2087,521)$,\allowbreak\ $(2136,31)$,\allowbreak\ $(2164,1221)$,\allowbreak\ $(2213,1074)$,\allowbreak\ $(2283,962)$,\allowbreak\ $(2311,780)$,\allowbreak\ $(2332,472)$,\allowbreak\ $(2360,633)$,\allowbreak\ $(2381,2383)$\\
\addlinespace[3pt]
11 & 1 & $0$ & $(0,2)$,\allowbreak\ $(3,1)$,\allowbreak\ $(3,2)$,\allowbreak\ $(3,8)$,\allowbreak\ $(7,0)$,\allowbreak\ $(7,5)$,\allowbreak\ $(7,6)$,\allowbreak\ $(8,1)$,\allowbreak\ $(8,3)$,\allowbreak\ $(8,7)$,\allowbreak\ $(9,8)$,\allowbreak\ $(10,6)$\\
\addlinespace[3pt]
13 & 1 & $0$ & $(0,7)$\\
13 & 1 & $1$ & $(6,0)$,\allowbreak\ $(7,10)$\\
13 & 2 & $0$ & $(9,136)$,\allowbreak\ $(22,136)$,\allowbreak\ $(35,136)$,\allowbreak\ $(48,136)$,\allowbreak\ $(61,136)$,\allowbreak\ $(74,136)$,\allowbreak\ $(87,136)$,\allowbreak\ $(100,136)$,\allowbreak\ $(113,136)$,\allowbreak\ $(126,136)$,\allowbreak\ $(139,136)$,\allowbreak\ $(152,136)$,\allowbreak\ $(165,136)$\\
\addlinespace[3pt]
\end{longtable}
\endgroup
\end{document}
